{{دیوان
| عنوان = مَیں اور تُو
| ماخذ = ویکی ماخذ
| ماخذ_لنک = https://ur.wikisource.org/wiki/%D9%85%DB%8C%DA%BA_%D8%A7%D9%88%D8%B1_%D8%AA%D9%88%D9%8F
| مقابلہ = اقبال سائبر لائبریری (اقبال اکادمی پاکستان)
| مقابلہ_لنک = https://www.iqbalcyberlibrary.net
}}
<poem>
مذاقِ دید سے ناآشنا نظر ہے مری
تری نگاہ ہے فطرت کی راز داں، پھر کیا

رہینِ شکوۂ ایّام ہے زبان مری
تری مراد پہ ہے دَورِ آسماں، پھر کیا

رکھا مجھے چمن آوارہ مثلِ موجِ نسیم
عطا فلک نے کِیا تجھ کو آشیاں، پھر کیا

فزوں ہے سُود سے سرمایۂ حیات ترا
مرے نصیب میں ہے کاوشِ زیاں، پھر کیا

ہوا میں تیرتے پھرتے ہیں تیرے طیّارے
مرا جہاز ہے محرومِ بادباں، پھر کیا

قوی شدیم چہ شد، ناتواں شدیم چہ شد
چنیں شدیم چہ شد یا چناں شدیم چہ شد

بہیچ گونہ دریں گلستاں قرارے نیست
توگر بہار شدی، ما خزاں شدیم، چہ شد
</poem>
